\section{Sobolev Spaces}\label{chap:sobolev-spaces}

The analysis of partial differential equations naturally involves function spaces which are defined not only in terms of properties of the functions themselves, but also of their derivatives. In Problem~\ref{ex:bounded-continuous-derivatives-banach}, for example, we introduced the Banach space $C_b^k(\overline\Omega)$. Unfortunately, these spaces turn out to be rather unsuitable for analyzing PDEs. For example, if $\Delta u=f$ with $f$ continuous, it is in general not true that $u$ is $C^2$. Sobolev spaces turn out to be much more useful. They are defined in an analogous fashion as $C_b^k$, but with $L^p$ taking over the role of the continuous bounded functions. Throughout this section, $\Omega$ will be an open, locally Jordan measurable set in $\mathbb{R}^m$.

This chapter is intended to give the reader a brief overview of Sobolev spaces. In order to arrange the chapter to make it easier for the reader to get the main ideas before jumping into the proofs we have arranged the material as follows.

\begin{itemize}
\item We begin with the basic definitions of the ``positive'' Sobolev spaces, examples of typical functions contained in (and excluded from) these spaces, and basic properties of the spaces.
\item We then have have a section on various important properties of Sobolev spaces. We address the following questions.
\begin{enumerate}
\item If $\Omega$ is all of $\mathbb{R}^m$ we can define the Fourier transform of a function in a Sobolev space. What do we know about the transform of a function in a particular space?
\item What imbedding relations exist between Sobolev spaces with different values of $p$? How large must $k$ be so that $W^{k,p}(\Omega)$ consists of continuous functions?
\item Compactness theorems: Recall the Arzela-Ascoli theorem which implies that a sequence of uniformly bounded continuous functions with uniformly bounded derivatives has a convergent subsequence. Can we replace uniform bounds on derivatives by $L^p$-bounds?
\item Trace theorems: Since $\mathcal{D}(\Omega)$ is dense in $L^p$, it is not meaningful to talk about boundary values of arbitrary $L^p$-functions. However, $\mathcal{D}(\Omega)$ is generally not dense in $W^{k,p}(\Omega)$. What can one say about boundary values of functions in Sobolev spaces? How can $W_0^{k,p}(\Omega)$ be characterized in terms of the behavior at the boundary?
\end{enumerate}
\item We then have a section on the dual spaces of Sobolev spaces, the so-called ``negative'' spaces.
\item Finally, we conclude with a section of technical results. These results are interesting in their own right, but are included here primarily to make it possible to prove some of the main results of the previous sections. The technical questions addressed are as follows.
\begin{enumerate}
\item Density theorems: Are functions in $C^k(\overline\Omega)$ dense in $W^{k,p}(\Omega)$? If $\Omega$ is unbounded, are functions of bounded support dense?
\item Coordinate transformations: How do functions in Sobolev spaces change under transformations to different domains?
\item Extension theorems: If $f\in W^{k,p}(\Omega)$, does there exist $F\in W^{k,p}(\mathbb{R}^m)$ such that $F|_{\Omega}=f$?
\end{enumerate}
\end{itemize}

\subsection{Basic Definitions}\label{sec:sobolev-basic-definitions}

\begin{Definition}
\label{def:sobolev-positive-integer-spaces}
Let $k$ be a non-negative integer and let $1\leq p\leq\infty$. Then we define $W^{k,p}(\Omega)$ to be the set of all distributions $u\in L^p(\Omega)$ such that $D^\alpha u\in L^p(\Omega)$ for $|\alpha|\leq k$. In $W^{k,p}(\Omega)$, we define a norm by
\begin{equation}\label{eq:sobolev-positive-integer-norm-definition}
\begin{aligned}
\|u\|_{k,p}^p&:=\sum_{|\alpha|\leq k}\|D^\alpha u\|_p^p,\quad p<\infty,\\
\|u\|_{k,\infty}&:=\max_{|\alpha|\leq k}\|D^\alpha u\|_{\infty},
\end{aligned}
\end{equation}
and for $p=2$, we define an inner product by
\begin{equation}\label{eq:sobolev-hilbert-inner-product-definition}
(u,v)_k:=\sum_{|\alpha|\leq k}\int_{\Omega}D^\alpha\overline{u(\bm{x})}D^\alpha v(\bm{x})\,\mathrm{d}\bm{x}.
\end{equation}
For $W^{k,2}(\Omega)$ we also use the notation $H^k(\Omega)$.
\end{Definition}
